%% SCRAP: experiments/campaigns/l8_attractor_map/results_20251209_145907/analysis_output/report/mathematical_framework
%% SOURCE: docs/working/experiments/campaigns/l8_attractor_map/results_20251209_145907/analysis_output/report/mathematical_framework.md
%% STATUS: HISTORICAL
%% FITS: experiments/ch-l8-map, ssrn/ch-math
%% EDITORIAL: lifted — prose rewritten to press voice

\section{L8 Attractor: Mathematical Framework for Heartbeat Dynamics}

The following equations characterize the temporal behavior of the StarForth
heartbeat observed in the 9 December 2025 run. The framework uses
quantum-thermodynamic analogs as modeling tools; the equations describe
measured phenomena, not claims about quantum physical mechanisms.

\subsection{Ground State (Equilibrium Frequency)}

The equilibrium angular frequency $\omega_0$ was measured independently
for each workload:

\begin{center}
\begin{tabular}{ll}
\toprule
Workload    & $\omega_0$ (Hz) \\
\midrule
DIVERSE     & $13.640 \pm 0.916$ \\
OMNI        & $13.430 \pm 0.794$ \\
STABLE      & $13.569 \pm 0.504$ \\
TEMPORAL    & $13.930 \pm 1.237$ \\
TRANSITION  & $13.731 \pm 1.978$ \\
VOLATILE    & $13.450 \pm 0.949$ \\
\bottomrule
\end{tabular}
\end{center}

The CV across all six workloads is 1.3\%, indicating a near-universal
ground-state frequency of approximately $\omega_0 \approx 13.5$\,Hz.

\subsection{Damped Harmonic Oscillator (Convergence Model)}

Convergence trajectories fit a damped harmonic oscillator:
\[
  \omega(t) = \omega_0 + A\,e^{-\gamma t}\cos(\Omega t + \varphi)
\]

Fitted parameters for selected workloads:

\begin{center}
\begin{tabular}{llll}
\toprule
Workload & $\gamma$ (tick$^{-1}$) & $\Omega$ (rad/tick) & $\tau = 1/\gamma$ (ticks) \\
\midrule
DIVERSE & 0.725 & 1.413 (period 4.45 ticks) & 1.4 \\
OMNI    & 0.045 & 0.450 (period 13.96 ticks) & 22.0 \\
STABLE  & 0.045 & 0.245 (period 25.67 ticks) & 22.4 \\
\bottomrule
\end{tabular}
\end{center}

\subsection{Boltzmann Distribution (Thermal Fluctuation Model)}

Using execution frequency variance as a proxy for thermal energy,
frequency fluctuations around $\omega_0$ are modeled by:
\[
  P(\omega) = \frac{1}{Z}\exp\!\left(\frac{-(\omega-\omega_0)^2}{k_B T}\right)
\]

Effective temperatures $T_\text{eff}$:

\begin{center}
\begin{tabular}{lll}
\toprule
Workload    & $k_B T$ (Hz$^2$) & $T_\text{eff}$ (Hz) \\
\midrule
STABLE      & 4.732 & 2.175 \\
VOLATILE    & 5.484 & 2.342 \\
OMNI        & 5.605 & 2.367 \\
TEMPORAL    & 6.678 & 2.584 \\
TRANSITION  & 7.240 & 2.691 \\
DIVERSE     & 7.483 & 2.735 \\
\bottomrule
\end{tabular}
\end{center}

\subsection{Entropy Production Rate}

Entropy production was estimated as:
\[
  \frac{dS}{dt} = \sum_i \frac{\Delta H_i}{T_i}
\]

Measured rates ranged from $3.8 \times 10^{-5}$ to $4.7 \times 10^{-5}$
(heat units per Hz per tick) across the six workloads.

\subsection{Frequency--Time Uncertainty}

An analog of the Heisenberg uncertainty relation,
$\Delta\omega \cdot \Delta t \geq \hbar/2$,
was observed in the data. Measured uncertainty products:

\begin{center}
\begin{tabular}{ll}
\toprule
Workload    & $\Delta\omega \cdot \Delta t$ (Hz$\cdot$s) \\
\midrule
STABLE      & $3.0 \times 10^{-5}$ \\
VOLATILE    & $4.0 \times 10^{-5}$ \\
OMNI        & $4.8 \times 10^{-5}$ \\
TEMPORAL    & $6.3 \times 10^{-5}$ \\
DIVERSE     & $8.9 \times 10^{-5}$ \\
TRANSITION  & $15.2 \times 10^{-5}$ \\
\bottomrule
\end{tabular}
\end{center}

TRANSITION, as the phase-change workload, exhibits the largest
uncertainty product — consistent with its high variability and the
difficulty of simultaneously characterizing its instantaneous frequency
and timing.

\subsection{Spectral Decomposition}

System behavior is modeled as a superposition of normal modes:
\[
  \omega(t) = \omega_0 + \sum_n A_n \cos(\omega_n t + \varphi_n)
\]

Each workload produces a distinct eigenmode spectrum. These spectral
signatures provide a basis for workload fingerprinting without explicit
runtime classification.

%% PATENT: spectroscopic workload fingerprinting via heartbeat eigenmode
%% spectra is patent-adjacent. Do not draft claims here.
